% ============================================================================
% APPENDIX LA: LIT-LAIBSON - Hyperbolic Discounting & Present Bias
% ============================================================================
% Category: Literature Integration (LIT-)
% Language: English
% EBF Integration: Time preferences, β-δ model, and commitment
% ============================================================================

\chapter*{Appendix UNMAPPED_LA: David Laibson Research Integration}
\addcontentsline{toc}{chapter}{Appendix UNMAPPED_LA: LIT-LAIBSON - Hyperbolic Discounting \& Present Bias}
\label{app:laibson}

% ============================================================================
% HEADER BLOCK
% ============================================================================
\begin{tcolorbox}[colback=white,colframe=black!75!white,title=Appendix Metadata]
\textbf{Appendix:} LA -- LIT-LAIBSON \\
\textbf{Category:} Literature Integration (LIT-) \\
\textbf{Version:} 1.0 (January 2026) \\
\textbf{Dependencies:} CORE-WHEN, LIT-RABIN, LIT-THALER, DOMAIN-FINANCE \\
\textbf{Status:} Complete
\end{tcolorbox}

% ============================================================================
% CROSS-REFERENCE MAP
% ============================================================================
\begin{tcolorbox}[colback=blue!5!white,colframe=blue!75!black,title=Cross-Reference Map]
\textbf{Builds On:}
\begin{itemize}
    \item Appendix UNMAPPED_V (CORE-WHEN): Temporal context $\Psi_t$
    \item Appendix U (LIT-KT): Prospect theory foundations
    \item Appendix R (LIT-THALER): Mental accounting, nudges
    \item Appendix UNMAPPED_RA (LIT-RABIN): Reference-dependent preferences
\end{itemize}

\textbf{Extends To:}
\begin{itemize}
    \item Appendix AD (DOMAIN-FINANCE): Savings and retirement
    \item Appendix AE (DOMAIN-HEALTH): Self-control in health behaviors
    \item Appendix S (PREDICT-BEHAVIOR): Behavioral prediction
    \item Appendix UNMAPPED_BBB (CORE-WHERE): Parameter repository
\end{itemize}

\textbf{Methods Connection:}
\begin{itemize}
    \item Appendix E (METHOD-E): Field experiments
    \item Appendix UNMAPPED_LLM (METHOD-LLMMC): Parameter estimation
\end{itemize}
\end{tcolorbox}

% ============================================================================
% CHAPTER LINKAGE
% ============================================================================
\begin{tcolorbox}[colback=green!5!white,colframe=green!75!black,title=Chapter Linkage]
\textbf{Primary:} Chapter 11 (Time Preferences) \\
\textbf{Secondary:} Chapter 6 (Utility), Chapter 14 (Financial Applications)
\end{tcolorbox}

% ============================================================================
% ABSTRACT
% ============================================================================
\begin{abstract}
This appendix integrates David Laibson's foundational research on hyperbolic discounting and present bias into the Evidence-Based Framework. Laibson's $\beta$-$\delta$ model revolutionized how economists model intertemporal choice by capturing the systematic tendency to overweight immediate rewards. His work on default effects in 401(k) plans demonstrated how choice architecture can leverage present bias to improve welfare. This appendix provides the core parameters for CORE-WHEN and establishes the theoretical foundation for commitment devices and paternalistic interventions.
\end{abstract}

% ============================================================================
% QUICK REFERENCE
% ============================================================================
\begin{tcolorbox}[colback=gray!10!white,colframe=gray!50!black,title=Quick Reference]
Key terms defined in this appendix (see also Appendix G for complete Glossary):
\begin{description}
    \item[$\beta$-$\delta$ Model] Quasi-hyperbolic discounting with present bias $\beta$ and long-run discount factor $\delta$
    \item[Present Bias] Tendency to overweight immediate outcomes ($\beta < 1$)
    \item[Naive vs. Sophisticated] Whether agents are aware of their future self-control problems
    \item[Commitment Device] Mechanism that restricts future choices to overcome present bias
    \item[Default Effect] Impact of pre-selected options on final choices
\end{description}
\end{tcolorbox}

% ============================================================================
% FUNDAMENTAL QUESTION
% ============================================================================
% -----------------------------------------------------------------------------
% APPENDIX SCOPE BOX
% -----------------------------------------------------------------------------

\begin{tcolorbox}[
    colback=orange!5!white,
    colframe=orange!75!black,
    title={\textbf{Appendix Scope: Ziel / In-Scope / Out-of-Scope / Constraints}},
    fonttitle=\bfseries
]

\textbf{Ziel (Objective):}
\begin{itemize}[nosep]
    \item Why do people systematically fail to act in their own long-term interest, and how can institutions help them overcome this tendency?
\end{itemize}

\textbf{In-Scope:}
\begin{itemize}[nosep]
    \item Fundamental Question
    \item Theoretical Framework
    \item Core Axioms
    \item Key Empirical Results
    \item EBF Integration
\end{itemize}

\textbf{Out-of-Scope (delegiert an andere Appendices):}
\begin{itemize}[nosep]
    \item CORE-WHEN $\rightarrow$ Appendix UNMAPPED_V
    \item LIT-KT $\rightarrow$ Appendix U
    \item LIT-THALER $\rightarrow$ Appendix R
    \item LIT-RABIN $\rightarrow$ Appendix UNMAPPED_RA
    \item DOMAIN-FINANCE $\rightarrow$ Appendix AD
\end{itemize}

\textbf{Constraints (Anwendungsgrenzen):}
\begin{itemize}[nosep]
    \item [TODO: Define when this appendix does NOT apply]
    \item [TODO: Define required prerequisites]
    \item [TODO: Define known limitations]
\end{itemize}

\textbf{Lieferobjekte (Deliverables):}
\begin{itemize}[nosep]
    \item 5 Table(s)
    \item 8 Equation(s)
    \item 6 Axiom(s)
    \item Worked Example(s)
\end{itemize}

\end{tcolorbox}


\section{Fundamental Question}

\begin{tcolorbox}[colback=yellow!10!white,colframe=orange!75!black,title=Central Research Question]
\textbf{Why do people systematically fail to act in their own long-term interest, and how can institutions help them overcome this tendency?}
\end{tcolorbox}

This question addresses CORE-WHEN ($\Psi_t$): How does the temporal dimension of context affect behavior?

% ============================================================================
% THEORETICAL FRAMEWORK
% ============================================================================
\section{Theoretical Framework}

\subsection{The $\beta$-$\delta$ Model}

Laibson's quasi-hyperbolic discounting model provides a parsimonious representation of present-biased preferences:

\begin{equation}
U_t = u_t + \beta \sum_{\tau=1}^{T} \delta^\tau u_{t+\tau}
\end{equation}

Where:
\begin{itemize}
    \item $u_t$ = Instantaneous utility at time $t$
    \item $\beta \in (0,1]$ = Present bias parameter (immediate gratification weight)
    \item $\delta \in (0,1)$ = Long-run discount factor (standard exponential)
    \item $\beta = 1$ reduces to standard exponential discounting
\end{itemize}

\subsection{Key Insight: Dynamic Inconsistency}

The model captures preference reversals:

\begin{align}
\text{At } t=0: \quad & \beta\delta^2 u_2 > \beta\delta u_1 \implies \text{prefer delayed} \\
\text{At } t=1: \quad & \beta\delta u_2 < u_1 \implies \text{prefer immediate}
\end{align}

This explains why people plan to save/exercise/diet but fail to follow through.

\subsection{Sophistication vs. Naivete}

\begin{table}[h]
\centering
\begin{tabular}{lcc}
\hline
\textbf{Type} & \textbf{Aware of $\beta < 1$?} & \textbf{Behavior} \\
\hline
Sophisticated & Yes & Seeks commitment \\
Naive & No & Perpetually procrastinates \\
Partially Naive & Underestimates & Insufficient commitment \\
\hline
\end{tabular}
\caption{Agent types in the $\beta$-$\delta$ framework}
\end{table}

% ============================================================================
% AXIOMS
% ============================================================================
\section{Core Axioms}

\begin{tcolorbox}[colback=red!5!white,colframe=red!75!black,title=Axiom UNMAPPED_LA-1: Present Bias Universality]
\textbf{Present bias ($\beta < 1$) is a robust feature of human decision-making across domains.}

Empirical estimate:
\begin{equation}
\beta \approx 0.70 \quad (SE = 0.05)
\end{equation}

Meaning: Future utility is discounted by an additional 30\% beyond standard exponential discounting.

\textit{Evidence:} Meta-analysis across savings, health, and labor domains.
\end{tcolorbox}

\begin{tcolorbox}[colback=red!5!white,colframe=red!75!black,title=Axiom UNMAPPED_LA-2: Default Power]
\textbf{Defaults have large effects on behavior because changing requires immediate effort for delayed benefit.}

\begin{equation}
P(\text{enroll} | \text{opt-in}) \approx 0.40 \quad \text{vs.} \quad P(\text{enroll} | \text{opt-out}) \approx 0.90
\end{equation}

The 50 percentage point gap persists even for financially consequential decisions.

\textit{Evidence:} Choi, Laibson, Madrian \& Metrick (2004), $N > 50,000$ employees.
\end{tcolorbox}

\begin{tcolorbox}[colback=red!5!white,colframe=red!75!black,title=Axiom UNMAPPED_LA-3: Shrouded Attribute Exploitation]
\textbf{Firms strategically hide costs that require future attention to discover.}

\begin{equation}
\text{Shrouding Equilibrium:} \quad c_{visible} \downarrow, \quad c_{shrouded} \uparrow, \quad \pi_{firm} \uparrow
\end{equation}

Present-biased consumers fail to account for shrouded fees (e.g., printer cartridges, bank fees, hotel surcharges).

\textit{Evidence:} Gabaix \& Laibson (2006), QJE.
\end{tcolorbox}

\begin{tcolorbox}[colback=red!5!white,colframe=red!75!black,title=Axiom UNMAPPED_LA-4: Commitment Demand]
\textbf{Sophisticated present-biased agents value commitment devices.}

\begin{equation}
WTP_{commitment} > 0 \iff \text{Agent is sophisticated} \land \beta < 1
\end{equation}

However, take-up is lower than welfare-optimal due to partial naivete.

\textit{Evidence:} Laibson (2015), behavioral mechanisms review.
\end{tcolorbox}

\begin{tcolorbox}[colback=red!5!white,colframe=red!75!black,title=Axiom UNMAPPED_LA-5: Age-Cognition Interaction]
\textbf{Financial decision quality follows an inverted-U shape with age.}

\begin{equation}
\text{Decision Quality}(age) = \alpha \cdot \text{Experience}(age) - \gamma \cdot \text{Cognitive Decline}(age)
\end{equation}

Peak performance: Age $\approx 53$ years.

\textit{Evidence:} Agarwal, Driscoll, Gabaix \& Laibson (2009), Brookings.
\end{tcolorbox}

\begin{tcolorbox}[colback=red!5!white,colframe=red!75!black,title=Axiom UNMAPPED_LA-6: Projection Bias]
\textbf{People systematically underestimate how much their preferences will change.}

\begin{equation}
\hat{u}_{t+1}(x | s_{t+1}) = (1-\alpha) \cdot u(x | s_{t+1}) + \alpha \cdot u(x | s_t)
\end{equation}

Where $\alpha \approx 0.4$ is the projection bias parameter.

\textit{Evidence:} Loewenstein, O'Donoghue \& Laibson (2003), QJE.
\end{tcolorbox}

% ============================================================================
% KEY EMPIRICAL RESULTS
% ============================================================================
\section{Key Empirical Results}

\subsection{The 401(k) Revolution}

\begin{table}[h]
\centering
\begin{tabular}{lccc}
\hline
\textbf{Intervention} & \textbf{Participation} & \textbf{Effect Size} & \textbf{Persistence} \\
\hline
Opt-in (baseline) & 37\% & -- & -- \\
Opt-out (auto-enroll) & 86\% & +49pp & 3+ years \\
Auto-escalation & +3\%/year & Cumulative & Permanent \\
\hline
\end{tabular}
\caption{Default effects in retirement savings}
\end{table}

\textbf{Key insight:} Defaults work not through information but through present bias---changing requires immediate action.

\subsection{The \$100 Bills on the Sidewalk}

Choi, Laibson \& Madrian (2011) documented:
\begin{itemize}
    \item 36\% of employees don't get full employer match
    \item Median ``free money'' left: \$507/year
    \item Financially literate employees also undercontribute
\end{itemize}

\textbf{Implication:} Present bias dominates even obvious financial gains.

\subsection{Parameter Estimates Across Studies}

\begin{table}[h]
\centering
\begin{tabular}{lccc}
\hline
\textbf{Domain} & \textbf{$\beta$} & \textbf{$\delta$ (annual)} & \textbf{Source} \\
\hline
Savings & 0.70 & 0.96 & Laibson et al. (1998) \\
Health & 0.65 & 0.95 & Various \\
Labor & 0.75 & 0.97 & DellaVigna \& Malmendier \\
\textbf{Consensus} & \textbf{0.70} & \textbf{0.96} & Meta-analysis \\
\hline
\end{tabular}
\caption{Present bias parameters across domains}
\end{table}

% ============================================================================
% EBF INTEGRATION
% ============================================================================
\section{EBF Integration}

\subsection{CORE-WHEN Specification}

Laibson's work provides the foundational parameters for temporal context:

\begin{equation}
\Psi_t = \begin{pmatrix}
\beta & \text{(present bias)} \\
\delta & \text{(long-run discount)} \\
\alpha_{proj} & \text{(projection bias)} \\
\sigma_{naive} & \text{(naivete degree)}
\end{pmatrix}
\end{equation}

\begin{table}[h]
\centering
\begin{tabular}{lccl}
\hline
\textbf{Parameter} & \textbf{Estimate} & \textbf{SE} & \textbf{Primary Source} \\
\hline
$\beta$ & 0.70 & 0.05 & Laibson (1997, 1998) \\
$\delta$ (annual) & 0.96 & 0.02 & Multiple studies \\
$\alpha_{proj}$ & 0.40 & 0.08 & Loewenstein et al. (2003) \\
Default effect & 0.50 & 0.05 & Choi et al. (2004) \\
\hline
\end{tabular}
\caption{EBF parameters from Laibson research}
\end{table}

\subsection{10C Framework Mapping}

\begin{table}[h]
\centering
\begin{tabular}{lll}
\hline
\textbf{CORE} & \textbf{Laibson Contribution} & \textbf{Key Paper} \\
\hline
WHO (AAA) & Self vs. future self as distinct agents & Harris \& Laibson (2001) \\
WHAT (C) & Immediate vs. delayed utility & Laibson (1997) \\
HOW (B) & Present-future self interaction & Ericson \& Laibson (2017) \\
WHEN (V) & $\beta$-$\delta$ temporal structure & All work \\
WHERE (BBB) & $\beta \approx 0.70$ empirical anchor & Meta-analysis \\
AWARE (AU) & Sophistication vs. naivete & Laibson (2015) \\
READY (AV) & Commitment device demand & Multiple \\
\hline
\end{tabular}
\caption{10C Framework mapping for Laibson research}
\end{table}

% ============================================================================
% WORKED EXAMPLE
% ============================================================================
\section{Worked Example: Designing Retirement Savings Nudges}

\begin{tcolorbox}[colback=blue!5!white,colframe=blue!75!black,title=Application: Auto-Enrollment Program Design]
\textbf{Context:} A company with 5,000 employees wants to improve retirement savings.

\textbf{Step 1: Baseline Assessment}
\begin{itemize}
    \item Current opt-in enrollment: 42\%
    \item Average contribution: 4\% of salary
    \item Employer match: 50\% up to 6\%
\end{itemize}

\textbf{Step 2: Apply $\beta$-$\delta$ Model}

Using Laibson parameters ($\beta = 0.70$):
\begin{itemize}
    \item Non-enrollees discount future benefits by additional 30\%
    \item Immediate ``hassle cost'' of enrollment blocks action
    \item Default removal eliminates present-bias barrier
\end{itemize}

\textbf{Step 3: Design Intervention}

\begin{enumerate}
    \item \textbf{Auto-enrollment} at 6\% (full match capture)
    \item \textbf{Auto-escalation} of 1\%/year up to 10\%
    \item \textbf{Simplification} of fund choices (target-date default)
\end{enumerate}

\textbf{Step 4: Predict Effects}

Using Choi et al. parameters:
\begin{align}
\text{New enrollment} &\approx 42\% + 50pp = 92\% \\
\text{New contribution} &\approx 6\% + (1\% \times 3 \text{ years}) = 9\% \\
\text{Additional savings/employee/year} &\approx \$2,500
\end{align}

\textbf{Step 5: Welfare Analysis}

Per Bernheim et al. (2015):
\begin{itemize}
    \item Welfare gain from auto-enrollment: \$0.50 per \$1 saved
    \item Opt-out rate: $<10\%$ (reveals preference alignment)
    \item Paternalism justified by present-bias correction
\end{itemize}
\end{tcolorbox}

% ============================================================================
% IMPLICATIONS
% ============================================================================
\section{Implications for Practice}

\subsection{For Choice Architects}

\begin{enumerate}
    \item \textbf{Set Smart Defaults:} Use opt-out, not opt-in
    \item \textbf{Leverage Auto-Escalation:} Gradual increases avoid immediate pain
    \item \textbf{Reduce Immediate Costs:} Simplify enrollment, pre-fill forms
    \item \textbf{Provide Commitment Options:} Allow voluntary restriction of future choices
\end{enumerate}

\subsection{For Policy}

\begin{enumerate}
    \item \textbf{Mandate Disclosure of Shrouded Fees:} Counter firm exploitation
    \item \textbf{Regulate Default Settings:} Require welfare-improving defaults
    \item \textbf{Age-Appropriate Protections:} Stronger safeguards for elderly
    \item \textbf{Commitment Account Options:} Tax-advantaged commitment savings
\end{enumerate}

% ============================================================================
% SUMMARY
% ============================================================================

% === AUTO-GENERATED ENTRIES (2026-02-22) ===

%%%%%%%%%%%%%%%%%%%%%%%%%%%%%%%%%%%%%%%%%%%%%%%%%%%%%%%%%%%%%%%%%%%%%%%%%%%%%%%
\subsubsection{1998: Self-Control and Saving for Retirement}
%%%%%%%%%%%%%%%%%%%%%%%%%%%%%%%%%%%%%%%%%%%%%%%%%%%%%%%%%%%%%%%%%%%%%%%%%%%%%%%

\paragraph{Citation.}
Laibson, David and Repetto, Andrea and Tobacman, Jeremy (1998). ``Self-Control and Saving for Retirement.'' \textit{Brookings Papers on Economic Activity}, 1998.

\paragraph{10C Integration.}
\begin{itemize}[nosep]
  \item \textbf{Domain}: [To be specified]
  \item \textbf{use\_for}: LIT-LAIBSON, LIT-O
\end{itemize}


%%%%%%%%%%%%%%%%%%%%%%%%%%%%%%%%%%%%%%%%%%%%%%%%%%%%%%%%%%%%%%%%%%%%%%%%%%%%%%%
\subsubsection{2015: Why Dont Present-Biased Agents Make Commitments?}
%%%%%%%%%%%%%%%%%%%%%%%%%%%%%%%%%%%%%%%%%%%%%%%%%%%%%%%%%%%%%%%%%%%%%%%%%%%%%%%

\paragraph{Citation.}
Laibson, David (2015). ``Why Dont Present-Biased Agents Make Commitments?.'' \textit{American Economic Review Papers and Proceedings}, 105.

\paragraph{10C Integration.}
\begin{itemize}[nosep]
  \item \textbf{Domain}: [To be specified]
  \item \textbf{use\_for}: LIT-LAIBSON, LIT-O
\end{itemize}


%%%%%%%%%%%%%%%%%%%%%%%%%%%%%%%%%%%%%%%%%%%%%%%%%%%%%%%%%%%%%%%%%%%%%%%%%%%%%%%
\subsubsection{2018: Private Paternalism, the Commitment Puzzle, and Model-Free E...}
%%%%%%%%%%%%%%%%%%%%%%%%%%%%%%%%%%%%%%%%%%%%%%%%%%%%%%%%%%%%%%%%%%%%%%%%%%%%%%%

\paragraph{Citation.}
Laibson, David (2018). ``Private Paternalism, the Commitment Puzzle, and Model-Free Equilibrium.'' \textit{AEA Papers and Proceedings}, 108.

\paragraph{10C Integration.}
\begin{itemize}[nosep]
  \item \textbf{Domain}: [To be specified]
  \item \textbf{use\_for}: LIT-LAIBSON, LIT-O
\end{itemize}

% === END AUTO-GENERATED ENTRIES ===
\section{Summary}

\begin{tcolorbox}[colback=gray!10!white,colframe=gray!50!black,title=Key Takeaways]
\begin{enumerate}
    \item \textbf{$\beta \approx 0.70$:} People discount the future by 30\% beyond exponential
    \item \textbf{Defaults Matter:} 50pp enrollment gap from opt-in to opt-out
    \item \textbf{Shrouding Works:} Firms exploit present bias to hide costs
    \item \textbf{Age Peak at 53:} Financial decision quality rises then falls
    \item \textbf{Commitment Demand:} Sophisticated agents value self-restriction
    \item \textbf{Paternalism Justified:} When correcting present bias, not preferences
\end{enumerate}

\textbf{Core Insight:} Design for the present-biased human, not the rational agent.
\end{tcolorbox}

% ============================================================================
% GLOSSARY
% ============================================================================
\section{Glossary}

Key terms defined in this appendix (see also Appendix G for complete Glossary):

\begin{description}
    \item[$\beta$-$\delta$ Model] Two-parameter discounting: $\beta$ for present bias, $\delta$ for exponential
    \item[Present Bias] Overweighting of immediate outcomes ($\beta < 1$)
    \item[Dynamic Inconsistency] Preference reversals over time
    \item[Sophisticated] Agent aware of own present bias
    \item[Naive] Agent unaware of own present bias
    \item[Commitment Device] Mechanism restricting future choice set
    \item[Default Effect] Behavioral impact of pre-selected options
    \item[Shrouded Attribute] Hidden cost exploiting consumer inattention
    \item[Auto-Escalation] Automatic increase in savings rate over time
\end{description}

% ============================================================================
% CRITICAL FOUNDATIONS
% ============================================================================
\section{Critical Foundations}

\begin{enumerate}
    \item \textbf{Strotz (1955):} Original analysis of dynamic inconsistency
    \item \textbf{Phelps \& Pollak (1968):} Intergenerational discounting model
    \item \textbf{Thaler \& Shefrin (1981):} Planner-doer model of self-control
    \item \textbf{O'Donoghue \& Rabin (1999):} Sophistication and naivete
    \item \textbf{Frederick, Loewenstein \& O'Donoghue (2002):} Time discounting review
\end{enumerate}

% ============================================================================
% OPEN ISSUES
% ============================================================================
\section{Open Issues}

\begin{enumerate}
    \item \textbf{Heterogeneity:} What determines individual $\beta$ variation?
    \item \textbf{Malleability:} Can present bias be reduced through training?
    \item \textbf{Domain Specificity:} Does $\beta$ vary across decision types?
    \item \textbf{Optimal Paternalism:} When does correction become manipulation?
    \item \textbf{Long-term Default Effects:} Do auto-enrollees become engaged savers?
\end{enumerate}

% ============================================================================
% REFERENCES
% ============================================================================
\section{References}

See Master Bibliography (bcm2\_master\_references.bib) for complete citations.

\nocite{bcm_master}

\textbf{Primary Laibson Sources} (18 papers integrated):

\begin{itemize}
    \item Laibson, D. (1997). Golden Eggs and Hyperbolic Discounting. \textit{QJE}, 112(2), 443--478.
    \item Harris, C. \& Laibson, D. (2001). Dynamic Choices of Hyperbolic Consumers. \textit{Econometrica}, 69(4).
    \item Choi, J., Laibson, D., Madrian, B., \& Metrick, A. (2002). Defined Contribution Pensions. \textit{Tax Policy \& Economy}.
    \item Choi, J., Laibson, D., Madrian, B., \& Metrick, A. (2003). Optimal Defaults. \textit{AER P\&P}.
    \item Choi, J., Laibson, D., Madrian, B., \& Metrick, A. (2004). For Better or For Worse. \textit{NBER}.
    \item Gabaix, X. \& Laibson, D. (2006). Shrouded Attributes. \textit{QJE}, 121(2).
    \item Agarwal, S., Driscoll, J., Gabaix, X., \& Laibson, D. (2009). Age of Reason. \textit{Brookings}.
    \item Choi, J., Laibson, D., \& Madrian, B. (2011). \$100 Bills on the Sidewalk. \textit{REStat}.
    \item Bernheim, D., Fradkin, A., Popov, I., \& Laibson, D. (2015). Welfare Economics of Defaults. \textit{AER}.
    \item Ericson, K. \& Laibson, D. (2017). Intertemporal Choice. \textit{Handbook of Behavioral Economics}.
\end{itemize}

\end{document}
